\begin{example}
	Посмотрим на отображение, которое переводит полярные координаты в обычные:
	\[
		f(r, \phi) = (r\cos\phi, r\sin\phi),\ f \colon [0; +\infty) \times \R \to \R^2
	\]
	Запишем его якобиан:
	\[
		\pd{(x, y)}{(r, \phi)} = \Det{&{\cos\phi} &{-r\sin\phi} \\ &{\sin\phi} &{r\cos\phi}} = r
	\]
	Заметим, что даже если $r > 0$, то у данной функции не может быть обратной. Действительно, если $f(r, \phi) = (x, y)$, то и для $f(r, \phi + 2\pi n) = (x, y),\ n \in \Z$
\end{example}

\begin{definition}
	Диффеоморфизмом класса $C^k$  называется отображение множества $G\subset\R^n$ в $\R^n$ являющееся взаимно однозначным, непрерывно дифференцируемым $k$ раз (здесь это означает непрерывную дифференцируемость $k$ раз каждой координатной функции, а последнее мы уже определяли). Обратное к которому обладает теми же свойствами.
\end{definition}